\addcontentsline{toc}{chapter}{Abstrak}
\chapter*{Abstrak}

\noindent
\begin{center}
	\begin{minipage}{0.9\linewidth}
		\centering
		\textbf{\MakeUppercase{\titleBahasa}}
	\end{minipage}
\end{center}

\vspace{2em}

\noindent
\begin{center}
	\renewcommand{\arraystretch}{1.25}
	\begin{tabularx}{0.75\linewidth}{@{}l@{\ }l@{\enspace}X}
		Nama Mahasiswa / NRP &:& \theAuthor\ / \theAuthorId \\
		Departemen &:& \departmentNameHeadlessBahasa\ \facultyNameBahasaAbbreviated\ - \campusNameAbbreviated \\
		Dosen Pembimbing / NIP &:& \supervisor\ / \supervisorId
	\end{tabularx}
\end{center}

\vspace{1em}

\section*{Abstrak}

Dimensi metrik dan dimensi metrik sisi pada sebuah graf merupakan invarian fundamental dalam teori graf. Secara sederhana, dimensi metrik menyatakan banyaknya simpul paling sedikit yang diperlukan untuk membedakan seluruh simpul lain dalam graf, sedangkan dimensi metrik sisi menyatakan hal serupa tetapi untuk membedakan seluruh sisi. Invarian ini memiliki berbagai aplikasi, seperti navigasi jaringan, pelacakan sumber, robotika, dan sistem biologis.

Tugas akhir ini mengeksplorasi penggunaan \foreign{Satisfiability Modulo Theories} (SMT) sebagai pendekatan untuk menghitung dimensi metrik dan dimensi metrik sisi. Permasalahan diformulasikan secara deklaratif dalam SMT dan diselesaikan menggunakan \foreign{solver}. Representasi graf kanonik digunakan untuk menghapus redundansi akibat automorfisma, dan eksperimen dilakukan pada seluruh graf terhubung unik dengan 2 hingga 9 simpul, serta beberapa graf berukuran lebih besar.

\vfill

\noindent
\begin{minipage}[t][0.375\textheight]{\textwidth}
	\begin{tabularx}{\linewidth}{@{}l@{\enspace}X}
		\textbf{Kata Kunci:} & \textit{Dimensi Metrik, Dimensi Metrik Sisi, \foreign{Satisfiability Modulo Theories}, SMT \foreign{Solver}, \foreign{NP-hard}}
	\end{tabularx}
\end{minipage}
